I want to draft an outline for my PhD thesis at University of Zurich. Official Start Date PhD: 01.10.2023. I am doing it at the Institute of Neuroinformatics. I want to give it a direction of we looked at deep learning and identified issues. Then, we did some techniques to reduce these issues. We argue that some neuro inspired solution could overcome the current limitations.

Part 1: About the current issues of AI Systems. We could use the following papers:


A Comprehensive Survey of Agents for Computer Use: Foundations, Challenges, and Future Directions

Abstract: Background: Agents for computer use (ACUs) are systems that execute complex tasks on digital devices – such as personal
computers or mobile phones – given instructions in natural language. These agents automate tasks by controlling software
through low-level actions like mouse clicks and touchscreen gestures. However, despite rapid progress, ACUs are not yet
mature for everyday use.
Objectives: This survey examines the current state-of-the-art, identifies trends, and points out research gaps in the develop-
ment of practical ACUs. The goal is to provide a comprehensive review and analysis that helps advance general-purpose,
robust, and scalable agents for real-world computer use.
Methods: We introduce a multifaceted taxonomy of ACUs across three dimensions: (I) the domain perspective, characterizing
the contexts in which agents operate; (II) the interaction perspective, describing observation modalities (e.g., screenshots,
HTML) and action modalities (e.g., mouse, keyboard, code execution); and (III) the agent perspective, detailing how agents
perceive, reason, and learn. We review 87 original research papers about ACUs and 33 relevant datasets, covering both
foundation model-based and specialized approaches.
Results: Our taxonomy comprehensively structures state-of-the-art approaches and establishes the groundwork for guiding
future ACU research. We found that the field is transitioning from specialized agents toward foundation-model-based agents, a
shift from text to image-based observation space, and an increasing adoption of behavior cloning methodologies. Furthermore,
we identify six key research gaps: insufficient generalization, inefficient learning, limited planning, low task complexity in
benchmarks, non-standardized evaluation, and a disconnect between research and practical conditions.
Conclusions: To continue rapid improvements in the field, we recommend focusing on: (a) vision-based observations and
low-level control to enhance generalization; (b) adaptive learning beyond static prompting; (c) effective planning and reasoning
capabilities; (d) realistic, high-complexity benchmarks; (e) standardized evaluation criteria based on task success; and (f)
aligning agent design with real-world deployment constraints. Collectively, our findings and proposed directions help develop
more general-purpose agents for everyday digital tasks.

My contribution: Shared First Author
Idea: Use this to say we reviewed one very promising field as an example. We found several issues.

This should be the main paper of this first part. Then, we have some papers to kind of strengthen the argument. But for these two papers, I am not a first authors and just did minor contributions. So maybe rather a half a page summary? Or what is the best way to deal with that?

So you want your private LLM at home? A survey and benchmark of methods for efficient GPTs

Abstract: At least since the introduction of ChatGPT, the
abilities of generative large language models (LLMs), sometimes
called GPTs, are at the center of the attention of AI researchers,
entrepreneurs, and others. However, for many applications, it
is not possible to call an existing LLM service via an API due
to data protection concerns or when no task-appropriate LLM
exists. On the other hand, deploying or training a private LLM is
often prohibitively computationally expensive. In this paper, we
give an overview of the most important recent methodologies that
help reduce the computational footprint of LLMs. We further
present extensive benchmarks for seven methods from two of
the most important areas of recent progress: model quantization
and low-rank adapters, showcasing how it is possible to leverage
state-of-the-art LLMs with limited resources. Our benchmarks
include resource consumption metrics (e.g. GPU memory usage),
a state-of-the-art quantitative performance evaluation as well as
a qualitative performance study conducted by eight individual
human raters. Our evaluations show that quantization has a
profound effect on GPU memory requirements. However, we also
show that these quantization methods, contrary to how they are
advertised, cause a noticeable loss in text quality. We further
show that low-rank adapters allow effective model fine-tuning
with moderate compute resources. For methods that require less
than 16 GB of GPU memory, we provide easy-to-use Jupyter
notebooks that allow anyone to deploy and fine-tune state-of-the-
art LLMs on the Google Colab free tier within minutes without
any prior experience or infrastructure.

My contribution: Evaluation of LLMs

and

COGITAO: A Visual Reasoning Framework To Study Compositionality & Generalization

Abstract: The ability to compose learned concepts and ap-
ply them in novel settings is key to human intelli-
gence, but remains a key challenge in state-of-the-
art machine learning models.
To address this issue, we introduce COGITAO, a
modulable data-generation framework to evaluate
compositional and systematic generalization in
object-centric domains.
Drawing inspiration from ARC-AGI’s environ-
ment and problem-setting, COGITAO constructs
rule-based tasks to be solved by applying a set of
transformations to objects in grid-based environ-
ments.
It supports composition over a set of 28 interoper-
able transformations, at adjustable composition-
depth, along with extensive control over grid
parametrization and object properties.
This flexibility enables the creation of millions of
unique task rules – surpassing existing datasets by
several orders of magnitude – across a broad range
of difficulties, while allowing virtually unlimited
sample generation per rule.
Alongside open-sourcing our flexible data-
generation framework, we release benchmark
datasets and provide baseline results with sev-
eral state-of-the-art architectures that incorporate
inductive biases well-suited for compositionality,
such as diffusion-based Transformers (LLaDA)
or recurrent Transformers with Adaptive Compu-
tation Time (Universal Transformer/PonderNet).
Despite strong in-domain performance, these
models consistently fail to generalize to novel


My contribution: I did the testing of MLDM on this dataset.


Ok, this is it of part 1. After this part, we have part 2

The following paper could be a main thing where we argue that backpropagation just learn statistical correlations.

The underlying structures of self-attention: symmetry, directionality, and emergent dynamics in Transformer training

Abstract: Self-attention is essential to Transformer architec-
tures, yet how information is embedded in the self-
attention matrices and how different objective func-
tions impact this process remains unclear. We present
a mathematical framework to analyze self-attention
matrices by deriving the structures governing their
weight updates. Using this framework, we demon-
strate that bidirectional training induces symmetry
in the weight matrices, while autoregressive train-
ing results in directionality and column dominance.
Our theoretical findings are validated across multi-
ple Transformer models — including ModernBERT,
GPT, LLaMA3, and Mistral — and input modalities
like text, vision, and audio. Finally, we apply these
insights by showing that symmetric initialization im-
proves the performance of encoder-only models on
language tasks. This mathematical analysis offers a
novel theoretical perspective on how information is
embedded through self-attention, thereby improving
the interpretability of Transformer models.

My contribution: Shared First Authors
My thoughts: Major part; I did most of the experiments but not much on the math side



After that part, we can have another part where we say there exist techniques to mitigate these flaws.

One Major part could be this one I am currently writing:

Multi-step prediction in learned world models compounds approximation errors exponentially: chaining a nonlinear function amplifies each step's error by its Lipschitz constant, causing long-horizon rollouts to diverge without bound. We show this failure is structural, a mathematical property of nonlinear function composition rather than a training deficiency, and address it by introducing \emph{Koopman-JEPA}, a self-supervised world model that lifts complex visual observations into a latent space where dynamics are globally linear. By optimizing a predictability-maximising JEPA objective instead of pixel-level reconstruction, the model makes this linear lifting tractable for high-dimensional visual inputs. Temporal rollouts thus reduce to simple matrix multiplications, mathematically bounding error accumulation to linear rather than exponential growth. We instantiate three variants with progressively stronger spectral constraints on the transition matrix, culminating in a Cayley parameterization that architecturally enforces strict stability.
Empirically, our Koopman variants reduce physical-state prediction error by up to 99\% on action-free tasks and 91\% on action-conditioned tasks at a 10-step horizon compared to a MLP-JEPA baseline, and maintain robust tracking on an action-free cartpole task where the baseline diverges catastrophically ($24.4 \pm 29.2$ vs.\ $0.250 \pm 0.012$). The linear inductive bias also improves sample efficiency by $\approx 16\times$ when using only $10$\% of the available training data. Finally, we provide a theoretical and numerical analysis establishing that only the Cayley parameterization satisfies the stabilizability precondition of discrete-time LQR, a finding confirmed empirically by a four-orders-of-magnitude difference in DARE gain norms versus unconstrained models.

Say we can impose hard constraints to account for current issues.

We can back it up with the following publications but those should only be minor things like half a page summary:

Hounsfield Unit Ranges as Inductive Bias for Intra-Clinical Learning of Data-Efficient CT Segmentation Models

Abstract: Automated tissue segmentation in medical imaging
plays a critical role in clinical AI-assisted decision-making, and
particularly in the assessment of body composition from CT
scans. However, acquiring data and annotations of sufficient
quality to train deep-learning models is expensive and time-
consuming. In this work, we propose a novel approach to
improve data efficiency and model accuracy by leveraging domain
knowledge about biologically relevant tissue-specific Hounsfield
unit (HU) ranges as an inductive bias for learning. Specifically, we
extend the input representation of deep learning-based segmen-
tation models with binary masks indicating potential tissue types,
where each binary mask is created from thresholds derived from
medical literature. Our method not only enhances segmentation
performance by up to 5 % for intramuscular adipose tissue but
surpasses the performance of the baseline model with 50 % of
the training data. Our easy-to-apply method thus improves data
efficiency and facilitates the development and use of segmentation
models in resource-constrained clinical settings.


My contribution: Shared First Authors, but did mostly the framing, not the implementation

and

Deep Retrieval at CheckThat! 2025: Identifying Scientific Papers from Implicit Social Media Mentions via Hybrid Retrieval and Re-Ranking

Abstract: We present the methodology and results of the Deep Retrieval team for subtask 4b of the CLEF CheckThat! 2025
competition, which focuses on retrieving relevant scientific literature for given social media posts. To address this
task, we propose a hybrid retrieval pipeline that combines lexical precision, semantic generalization, and deep
contextual re-ranking, enabling robust retrieval that bridges the informal-to-formal language gap. Specifically,
we combine BM25-based keyword matching with a FAISS vector store using a fine-tuned INF-Retriever-v1 model
for dense semantic retrieval. BM25 returns the top 30 candidates, and semantic search yields 100 candidates,
which are then merged and re-ranked via a large language model (LLM)-based cross-encoder.
Our approach achieves a mean reciprocal rank at 5 (MRR@5) of 76.46% on the development set and 66.43%
on the hidden test set, securing the 1st position on the development leaderboard and ranking 3rd on the test
leaderboard (out of 31 teams), with a relative performance gap of only 2 percentage points compared to the
top-ranked system. We achieve this strong performance by running open-source models locally and without
external training data, highlighting the effectiveness of a carefully designed and fine-tuned retrieval pipeline.

My contribution: First Authors

For the second one we could argue that another constraint is to look up knowledge during inference to prevent some of the issues.


Then, the next part is a change of paradigm. It is about different learning ideas could solve some of the issues


The Cooperative Network Architecture: Learning Structured Networks as Representation of Sensory Patterns

Abstract: We introduce the Cooperative Network Architecture (CNA), a model that represents sensory signals
using structured, recurrently connected networks of neurons, termed “nets.” Nets are dynamically
assembled from overlapping net fragments, which are learned based on statistical regularities in
sensory input. This architecture offers robustness to noise, deformation, and generalization to out-
of-distribution data, addressing challenges in current vision systems from a novel perspective. We
demonstrate that net fragments can be learned without supervision and flexibly recombined to encode
novel patterns, enabling figure completion and resilience to noise. Our findings establish CNA as a
promising paradigm for developing neural representations that integrate local feature processing with
global structure formation, providing a foundation for future research on invariant object recognition.

My contribution: First Authors
My thoughts: Major part


Then I still plan to get a paper using Deep Feedback Control (DFC). But this is not done yet. So the structure should be in way that this can be added here as a second work with different paradigm but that it also would work without it.


Given this, please help me first draft an outline of the thesis, a kind of concept on whose basis we then write the abstract.

